% Generated by scripts/assemble.py; corrections belong in source/corrections.json.
\examhead{198}{湘潭大学 2026 年数学分析试卷}{187}
% p188-b01; source PDF page 188
\blockanchor{p188-b01}
\question{1}(30 分) 计算题.\par
% p188-b02; source PDF page 188
\blockanchor{p188-b02}
(1) (10 分) 求极限  $\lim_{x\to0}\left(\cos x-\frac{x^{2}}{2}\right)^{\frac{1}{x^{2}}}$ .\par
% p188-b03; source PDF page 188
\blockanchor{p188-b03}
(2) (10 分) 求极限 $\lim_{x\to 0}x\left[\frac{1}{x}\right]$.\par
% p188-b04; source PDF page 188
\blockanchor{p188-b04}
(3) (10 分) 求积分 $\int_{0}^{\frac{\pi}{2}} \frac{\sin^{2} x}{\sin x + \cos x} \, \mathrm{d}x$.\par
% p188-b05; source PDF page 188
\blockanchor{p188-b05}
\question{2}(20 分) 计算题.\par
% p188-b06; source PDF page 188
\blockanchor{p188-b06}
(1) (10 分) 设 $y = \sqrt[3]{x - \tan x}$, 求 $\mathrm{d}y$ 及 $\mathrm{d}^2 y$.\par
% p188-b07; source PDF page 188
\blockanchor{p188-b07}
(2) (10 分) 设 $F(y) = \int_{a}^{b} f(x) |y - x| \, \mathrm{d}x$, 其中 $f(x)$ 可微, 求 $F''(y)$.\par
% p188-b08; source PDF page 188
\blockanchor{p188-b08}
\question{3}(15 分) 计算重积分 $\iiint_{V} \sqrt{1 - \frac{x^2}{a^2} - \frac{y^2}{b^2} - \frac{z^2}{c^2}} \, \mathrm{d}x \, \mathrm{d}y \, \mathrm{d}z$, 其中 $V = \left\{(x, y, z) \mid \frac{x^2}{a^2} + \frac{y^2}{b^2} + \frac{z^2}{c^2} \leq 1\right\}$.\par
% p188-b09; source PDF page 188
\blockanchor{p188-b09}
\question{4}(15 分) 计算曲面积分 $\iint_{\Sigma} \frac{x \, \mathrm{d}y \, \mathrm{d}z + y \, \mathrm{d}z \, \mathrm{d}x + z \, \mathrm{d}x \, \mathrm{d}y}{(x^2 + y^2 + z^2)^{\frac{3}{2}}}$, 其中 $\Sigma: x^2 + 2y^2 + 3z^2 = 1$, 取外侧.\par
% p188-b10; source PDF page 188
\blockanchor{p188-b10}
\question{5}(20 分) 讨论广义积分 $I = \int_{0}^{1} x^{p-1} \ln^{2} x \, \mathrm{d}x$ 的一致收敛性:\par
% p188-b11; source PDF page 188
\blockanchor{p188-b11}
(1) (10 分) $p \geq p_0 > 0$.\par
% p188-b12; source PDF page 188
\blockanchor{p188-b12}
(2) (10 分) p > 0.\par
% p188-b13; source PDF page 188
\blockanchor{p188-b13}
\question{6}(20 分) 设 $f(x)$ 在 $[a, b]$ 上连续.\par
% p188-b14; source PDF page 188
\blockanchor{p188-b14}
(1) (10 分) 证明: $f(x)$ 在 $[a, b]$ 上有界.\par
% p188-b15; source PDF page 188
\blockanchor{p188-b15}
(2) (10 分) 证明: $f(x)$ 在 $[a, b]$ 上一定有最大值和最小值.\par
% p188-b16; source PDF page 188
\blockanchor{p188-b16}
\question{7}(20 分) 设 $S_{n}(x) = n^{\alpha}xe^{-nx}$, 其中 $\alpha$ 是参数, 求 $\alpha$ 的取值范围, 使得函数列 $\{S_{n}(x)\}$ 在 $[0,1]$ 上:\par
% p188-b17; source PDF page 188
\blockanchor{p188-b17}
(1) (6 分) 一致收敛.\par
% p188-b18; source PDF page 188
\blockanchor{p188-b18}
(2) (6 分) 积分运算与极限运算可以交换, 即 $\lim_{n\to \infty}\int_0^1 S_n(x)\mathrm{d}x = \int_0^1\lim_{n\to \infty}S_n(x)\mathrm{d}x.$\par
% p188-b19; source PDF page 188
\blockanchor{p188-b19}
(3) (8 分) 求导运算与极限运算可以交换, 即对一切 $x \in [0,1]$, 成立 $\lim_{n \to \infty} \frac{\mathrm{d}}{\mathrm{d}x} S_n(x) = \frac{\mathrm{d}}{\mathrm{d}x} \lim_{n \to \infty} S_n(x)$.\par
% p188-b20; source PDF page 188
\blockanchor{p188-b20}
\question{8}(10 分) 求 $\cos x$ 的全部零点的倒数平方之和.\par
